
\subsubsection{3.1 Actionable Specificity Framework}

The proposed framework hypothesizes that feedback effectiveness for LLM code repair depends on Actionable Specificity (AS)---the degree to which a verification signal provides information that enables targeted bug localization and fix inference. AS is decomposed into three orthogonal dimensions:

\textbf{AS\_loc (Localization):} Whether the signal identifies the failure location. Binary: 1 if file:line information is present, 0 otherwise. Extraction pattern: \texttt{File ''([\^{}'']+)'', line (\textbackslash textbackslash d+)}.

\textbf{AS\_state (State Exposure):} How many variable values are visible at the failure point. Continuous: count of variable-value pairs. Extraction pattern: \texttt{(\textbackslash textbackslash w+)\textbackslash textbackslash s\textbackslash textit\{=\textbackslash textbackslash s\}([''''']?[\textbackslash textbackslash w\textbackslash textbackslash d\textbackslash .\textbackslash -\textbackslash [\textbackslash ]\{\}]+[''''']?)}.

\textbf{AS\_causal (Causal Context):} Execution trace depth between variable assignment and failure. Continuous: count of traceback frames. Extraction pattern: \texttt{\^{}\textbackslash textbackslash s+File ''([\^{}'']+)'', line (\textbackslash textbackslash d+), in (\textbackslash textbackslash w+)}.

These components are designed to be independently measurable from signal text using regex extraction.

\subsubsection{3.2 Signal Generation Conditions}

Six signal generation conditions were defined, varying in expected AS level:

\begin{table}[htbp]
\centering
\resizebox{\columnwidth}{!}{%
\begin{tabular}{lll}
\toprule
Condition & Description & Expected AS Level \\
\midrule
C1 & Full runtime trace & High (all components) \\
C2 & Truncated trace (last 3 frames) & Medium-High \\
C3 & Value-masked trace & Medium (no AS\_state) \\
C4 & Error message only & Low \\
C5 & Static analysis (placeholder) & Variable \\
C6 & Syntax error baseline & Minimal \\
\bottomrule
\end{tabular}
}
\end{table}

The conditions are ordered such that C1 $\geq$ C2 $\geq$ C3 $\geq$ C4 in expected AS, enabling ordered hypothesis testing.

\subsubsection{3.3 Design Rationale}

The simplest operationalization (regex extraction) was chosen to establish feasibility bounds before investing in more sophisticated methods (AST analysis, LLM-based extraction). This follows the principle of testing necessary conditions first: if AS components cannot be reliably extracted via simple patterns, the framework requires reformulation before testing predictive power.

\subsubsection{3.4 Existence Test Design (H-E1)}

Before testing whether AS predicts repair success, the existence hypothesis tested whether AS is measurable:

> \textbf{H-E1:} Under standard verification signal generation, AS\_loc, AS\_state, and AS\_causal can be independently computed for $\geq 95$\% of signals, and AS values show expected ordering across conditions ($C1\geq C2\geq C3\geq C4)$.

\textbf{Success Criteria:}
\begin{itemize}
\item Primary: Extraction rate $\geq 95$\% across all signals
\item Secondary: AS ordering $C1\geq C2\geq C3\geq C4$ preserved
\end{itemize}

